\chapter{Tecnologias e Ferramentas Utilizadas}
\label{sec:tecnologias}

Para o desenvolvimento da plataforma de monitorização e análise
da qualidade do ar em Beja, foi necessário selecionar um conjunto
de tecnologias que permitisse implementar as diferentes
funcionalidades do sistema, desde a recolha e tratamento dos dados
até ao seu armazenamento e posterior visualização.

A escolha das tecnologias teve em consideração os requisitos
funcionais da aplicação, a facilidade de integração entre os
diferentes componentes, a possibilidade de manutenção e a
adequação das ferramentas aos conhecimentos adquiridos ao longo
da licenciatura em Engenharia Informática.

A arquitetura adotada divide-se em três componentes principais:
o \textit{backend}, responsável pelo tratamento dos dados e
pela disponibilização da API REST; a base de dados, destinada
ao armazenamento estruturado das leituras; e o
\textit{frontend}, responsável pela apresentação e interação
com o utilizador.


\subsection{Linguagens de Programação e Bibliotecas}
\label{subsec:linguagens}

O desenvolvimento da aplicação envolveu a utilização de
diferentes linguagens de programação, frameworks e bibliotecas,
cada uma com responsabilidades específicas na arquitetura
do sistema.


\subsubsection{Python}

A linguagem Python foi utilizada no desenvolvimento do
\textit{backend}, nomeadamente na implementação do processo
ETL (\textit{Extract, Transform, Load}) e na gestão das
operações relacionadas com os dados ambientais.

A escolha desta linguagem deveu-se à sua sintaxe relativamente
simples, facilidade de manutenção e disponibilidade de
bibliotecas adequadas ao processamento de dados.

No processo ETL, foram utilizadas diferentes bibliotecas
da linguagem Python, destacando-se:

\begin{itemize}

    \item \textbf{csv}: responsável pela leitura dos ficheiros
    CSV e interpretação dos respetivos campos;

    \item \textbf{io}: utilizada para manipular o conteúdo
    dos ficheiros em memória;

    \item \textbf{datetime}: utilizada para interpretar,
    converter e validar os campos temporais;

    \item \textbf{zoneinfo}: responsável pela conversão
    dos timestamps para o fuso horário
    \texttt{Europe/Lisbon};

    \item \textbf{math}: utilizada na validação dos valores
    numéricos, nomeadamente para identificar valores
    não finitos;

    \item \textbf{sqlite3}: utilizada para estabelecer
    a comunicação entre Python e a base de dados SQLite.

\end{itemize}

Estas bibliotecas permitiram desenvolver um processo ETL
capaz de interpretar os ficheiros disponibilizados,
normalizar os atributos, validar a consistência das
leituras e preparar os dados para armazenamento.


\subsubsection{Flask}

Para implementar a comunicação entre o
\textit{frontend} e o \textit{backend}, foi utilizado
o framework Flask, desenvolvido para a linguagem Python.

O Flask permite criar aplicações web e disponibilizar
serviços através de uma API REST
(\textit{Representational State Transfer}).

No contexto deste projeto, a API permite consultar
os sensores existentes, obter leituras ambientais,
importar ficheiros CSV e eliminar dados armazenados.

Foram implementados os seguintes endpoints:

\begin{itemize}

    \item \texttt{GET /api/sensores}:
    devolve a lista de sensores registados;

    \item \texttt{GET /api/leituras}:
    permite consultar as leituras armazenadas,
    podendo filtrar os resultados por sensor;

    \item \texttt{POST /api/upload}:
    recebe ficheiros CSV, executa o processamento
    ETL e armazena os dados válidos;

    \item \texttt{DELETE /api/dados}:
    permite eliminar os dados registados
    na base de dados.

\end{itemize}

A comunicação entre os componentes é realizada
através de pedidos HTTP, sendo os resultados
disponibilizados em formato JSON
(\textit{JavaScript Object Notation}).

Foi também utilizada a biblioteca
\texttt{Flask-CORS}, permitindo configurar
as permissões de comunicação entre o servidor
Flask e o ambiente de desenvolvimento do Vue.js.


\subsubsection{Vue.js}

Para o desenvolvimento da interface web, foi utilizado
o framework JavaScript Vue.js, na sua terceira versão.

A escolha do Vue.js teve como principal objetivo
desenvolver uma interface interativa e organizada,
permitindo apresentar os dados ambientais de forma
dinâmica.

O Vue.js disponibiliza uma arquitetura baseada
em componentes, facilitando a separação das
responsabilidades e a reutilização de elementos
da interface.

A aplicação encontra-se organizada em diferentes
componentes, nomeadamente:

\begin{itemize}

    \item \textbf{TopHeader.vue}:
    apresenta o cabeçalho, permite selecionar
    sensores e importar ficheiros CSV;

    \item \textbf{TimeSlider.vue}:
    permite navegar temporalmente entre os
    diferentes registos;

    \item \textbf{AqiGaugeCard.vue}:
    apresenta visualmente o índice de qualidade
    do ar;

    \item \textbf{AqiMap.vue}:
    apresenta a localização geográfica dos
    sensores;

    \item \textbf{AqiEvolutionChart.vue}:
    permite visualizar a evolução histórica
    das concentrações dos poluentes;

    \item \textbf{PollutantList.vue}:
    apresenta as concentrações dos poluentes
    e os respetivos valores de referência.

\end{itemize}

A organização por componentes permite
desenvolver e modificar funcionalidades
específicas sem alterar diretamente
toda a estrutura da aplicação.

O ficheiro \texttt{App.vue} funciona como
componente principal, sendo responsável pela
gestão do estado da aplicação e pela
comunicação entre os diferentes componentes.


\subsubsection{Bibliotecas de Visualização e Comunicação}

Para complementar as funcionalidades do Vue.js,
foram utilizadas bibliotecas JavaScript específicas.

A biblioteca \textbf{Chart.js}, integrada através
de \texttt{vue-chartjs}, foi utilizada para
representar graficamente a evolução histórica
das concentrações de poluentes atmosféricos.

A utilização de gráficos de linhas permite
observar variações temporais e identificar
períodos com concentrações mais elevadas.

Para a representação geográfica dos sensores,
foi utilizada a biblioteca \textbf{Leaflet},
que permite integrar mapas interativos
na aplicação web.

Os sensores são representados através de
marcadores geográficos, cuja cor varia
de acordo com a classificação do índice
de qualidade do ar calculado pela aplicação.

Para a comunicação entre o Vue.js e a API REST
Flask, foi utilizada a biblioteca \textbf{Axios},
responsável pela realização de pedidos HTTP
e pelo processamento das respostas recebidas.

Por último, a biblioteca \textbf{Bootstrap}
foi utilizada para apoiar a organização visual
da interface, disponibilizando classes e
componentes que facilitam a construção
de um layout responsivo.


\subsection{Sistema de Gestão da Base de Dados}
\label{subsec:sqlite}

Para o armazenamento dos dados ambientais,
foi selecionado o sistema de gestão de base
de dados relacional SQLite.

O SQLite permite criar e gerir uma base de
dados através de um ficheiro local, não
exigindo a instalação ou configuração de
um servidor de base de dados independente.

Esta característica revelou-se adequada
ao contexto do projeto, uma vez que
simplifica a configuração do ambiente
de desenvolvimento e permite gerir
os dados de forma estruturada.

A base de dados foi organizada em três
tabelas principais:

\begin{itemize}

    \item \textbf{sensores}:
    armazena os identificadores dos sensores,
    nomes, coordenadas geográficas
    e respetivos metadados;

    \item \textbf{ficheiros}:
    regista os ficheiros CSV importados
    e identifica os sensores aos quais
    estão associados;

    \item \textbf{leituras}:
    armazena os registos ambientais,
    incluindo os timestamps, concentrações
    de poluentes, variáveis meteorológicas
    e coordenadas.

\end{itemize}

A utilização de chaves primárias e
estrangeiras permite estabelecer relações
entre as tabelas, assegurando a integridade
referencial dos dados.

Adicionalmente, foram definidas restrições
de unicidade e índices para facilitar
a organização e consulta das leituras.

A comunicação com o SQLite é realizada
através da biblioteca \texttt{sqlite3},
integrada na biblioteca padrão do Python.

O ficheiro \texttt{database.py} é responsável
pela criação da estrutura relacional,
estabelecimento das ligações e execução
das consultas necessárias ao funcionamento
da aplicação.


\subsection{Ambiente de Desenvolvimento e Controlo de Versões}
\label{subsec:ambiente}

O desenvolvimento do projeto foi realizado
no Visual Studio Code, um editor de código
que disponibiliza funcionalidades de apoio
ao desenvolvimento de aplicações em
diferentes linguagens de programação.

A sua utilização permitiu trabalhar com
os ficheiros Python, JavaScript, Vue.js,
HTML e CSS num único ambiente.

O Visual Studio Code disponibiliza ainda
um terminal integrado, utilizado para
executar os comandos necessários à
instalação de dependências, inicialização
do servidor Flask e execução da
aplicação Vue.js.

No desenvolvimento do \textit{frontend},
foi utilizada a ferramenta \textbf{Vite},
responsável pela disponibilização
do servidor de desenvolvimento e pela
compilação da aplicação Vue.js.

A gestão das dependências JavaScript foi
realizada através do \textbf{npm}
(\textit{Node Package Manager}), enquanto
as dependências Python foram instaladas
através do \textbf{pip}.

Para isolar as dependências do ambiente
Python, foi utilizado um ambiente virtual,
permitindo separar as bibliotecas do
projeto das restantes instalações
existentes no sistema.

Relativamente ao controlo de versões,
o Git constitui uma ferramenta adequada
para acompanhar a evolução do código,
registar alterações e facilitar a
recuperação de versões anteriores.

A utilização de um repositório GitHub
permite ainda centralizar o código-fonte
e manter um histórico do desenvolvimento,
caso seja adotado como plataforma de
alojamento do repositório.


\subsection{Justificação das Escolhas Tecnológicas}
\label{subsec:justificacao}

A seleção das tecnologias procurou
estabelecer um equilíbrio entre
facilidade de desenvolvimento,
organização do código, capacidade
de processamento e manutenção futura.

A utilização de Python permitiu
implementar o processo ETL através
de bibliotecas adequadas à leitura,
normalização e validação dos dados.

O framework Flask foi selecionado
pela sua simplicidade e facilidade
de integração com Python, permitindo
disponibilizar uma API REST sem
introduzir complexidade desnecessária
na arquitetura do sistema.

A utilização de SQLite revelou-se
adequada para a dimensão e os
objetivos iniciais do projeto,
permitindo armazenar os dados
de forma relacional e garantir
a consistência das relações
entre sensores, ficheiros e leituras.

No desenvolvimento do frontend,
o Vue.js permitiu construir uma
interface modular e reativa,
facilitando a atualização automática
dos elementos visuais quando
o utilizador seleciona diferentes
sensores ou períodos temporais.

Por sua vez, a integração de
Chart.js e Leaflet permitiu
representar os dados de forma
gráfica e geográfica, contribuindo
para uma interpretação mais
acessível da informação ambiental.

A separação entre frontend,
backend e base de dados contribui
para uma arquitetura mais organizada,
em que cada componente apresenta
responsabilidades específicas.

Esta abordagem facilita a manutenção
da aplicação e permite considerar
futuras evoluções, como a integração
de novas fontes de dados,
a implementação de funcionalidades
de análise adicionais ou a
substituição do sistema de
armazenamento, caso os requisitos
venham a justificar essa alteração.
